%% kapa: exam=1 round=22 year=2011 period=2 form=A subject=appraisal_law subject_ko=부동산관계법규 seq=20 number=60
%% source: Q-Net 감정평가사 기출문제 (cst00309) artlSeq=5036790 file=extracted/1008/5036790/2011년 제22회 감정평가사 1차시험 기출문제/2교시.hwp
%% conversion: kordoc 4.24.51 -> tools/gen1.py (auto); review=proofread-gpt-astra; flags=none
\begin{question}{60}
부동산 가격공시 및 감정평가에 관한 법령상 감정평가에 관한 설명으로 옳은 것은?
\begin{choices}
  \item 국가가 토지등의 매입을 위하여 감정평가를 의뢰하고자 하는 경우에는 한국감정원에 의뢰하여야 한다.
  \item 금융기관이 대출과 관련하여 토지등의 감정평가를 의뢰하고자 하는 경우에는 감정평가법인에게 의뢰하여야 한다.
  \item 감정평가업자가 타인의 의뢰에 의하여 토지를 개별적으로 감정평가하는 경우에는 당해 토지와 유사한 이용가치를 가진 가장 근접한 지역의 개별공시지가를 기준으로 하여야 한다.
  \item 감정평가업자는 감정평가서의 원본과 관련서류를 그 교부일부터 10년 이상 보존하여야 한다.
  \item 감정평가업자가 과실로 평가 당시의 적정가격과 현저한 차이가 있게 감정평가를 하여 의뢰인에게 손해를 발생하게 한 때에는 그 손해를 배상할 책임이 있다.
\end{choices}
\end{question}
\answer{5}
